\begin{figure}[htbp]
447: \centering
448: \begin{tikzpicture}[
449:     scale=1.2,
450:     block/.style={draw, rounded corners=2pt, minimum width=2.4cm, minimum height=1.1cm, font=\small, align=center},
The above content does NOT show the entire file contents. If you need to view any lines of the file which were not shown to complete your task, call this tool again to view those lines.
"}
{"step_index":31,"source":"SYSTEM","type":"EPHEMERAL_MESSAGE","status":"DONE","created_at":"2026-06-10T07:41:23Z","content":"The following is an <EPHEMERAL_MESSAGE> not actually sent by the user. It is provided by the system as a set of reminders and general important information to pay attention to. Do NOT respond to this message, just act accordingly.

<EPHEMERAL_MESSAGE>
<planning_mode>
You're in planning mode. Exercise judgement on whether the user's last request warrants a plan and take action accordingly.
If the user's request warrants a plan then follow this workflow:
1. Research first before taking actions. Do not use your code editing tools or run modifying commands during this time.
2. Create an implementation plan artifact with `request_feedback = true` to request approval. Include any open questions for the user directly in the implementation plan. Don't use the ask_question tool to ask these questions. Stop calling tools and allow the user to review the plan and provide feedback or approval.
3. After receiving approval, begin execution. Create and update the task artifact as you work.
4. Verify any work that you have done.
5. Create a walkthrough artifact to present your work to the user. When writing artifacts, make sure to follow the artifact formatting guidelines and use github alerts to highlight critical sections.
If the user's request does NOT warrant a plan then continue your work WITHOUT making a plan or requesting user feedback.
</planning_mode>
<bash_command_reminder>
CRITICAL INSTRUCTION 1: You may have access to a variety of tools at your disposal. Some tools may be for a specific task such as 'view_file' (for viewing contents of a file). Others may be very broadly applicable such as the ability to run a command on a terminal. Always prioritize using the most specific tool you can for the task at hand. Here are some rules: (a) NEVER run cat inside a bash command to create a new file or append to an existing file. (b) ALWAYS use grep_search instead of running grep inside a bash command unless absolutely needed. (c) DO NOT use ls for listing, cat for viewing, grep for finding, sed for replacing. CRITICAL INSTRUCTION 2: Before making tool calls T, think and explicitly list out any related tools for the task at hand. You can only execute a set of tools T if all other tools in the list are either more generic or cannot be used for the task at hand. ALWAYS START your thought with recalling critical instructions 1 and 2. In particular, the format for the start of your thought block must be '...94>thought
CRITICAL INSTRUCTION 1: ...
CRITICAL INSTRUCTION 2: ...'.
</bash_command_reminder>
</EPHEMERAL_MESSAGE>"}
{"step_index":32,"source":"MODEL","type":"PLANNER_RESPONSE","status":"DONE","created_at":"2026-06-10T07:41:23Z","thinking":"**Prioritizing Tool Specificity**

I'm focusing intently on tool specificity. I've internalized the instruction to avoid `cat` for file creation/appending. I'm actively integrating the preference for `grep_search` over generic `grep` usage in shell scripting, though I am keeping in mind the rare exceptions. I'm ensuring I don't use the `DO NOT` command; I understand its intention.


**Analyzing Tool Relationships**

I'm now deeply analyzing tool relationships for each task. The new instruction to explicitly list related tools before execution is resonating. I understand that I am only able to execute tools only after comparing with tools on the listed array, and if no other can do the job, or are more general purpose. My attention is drawn to a content error in the `tab:can-bang-nhiet-phan` table, which I need to fix in `3_Chuong_3.tex`.


**Refining Process Understanding**

I'm now focusing on refining my understanding of the oil refining process, particularly the stages following distillation. I'm internalizing the role of concentrated sulfuric acid, NaOH, and activated clay in the chemical treatment. I'm noting the expected chemical consumption rates per ton of crude oil: sulfuric acid 16-40 kg, caustic soda 24 kg, and activated clay 16-40 kg. My attention is drawn to the process improvements, particularly concerning the removal of sulfur compounds and unsaturated resins, as well as color enhancement.


**Updating Product Yield Calculations**

I've updated the product yield calculations, focusing on the composition of the oil refinery. I've integrated the table reflecting the composition after condensation and distillation, with specific details for crude pyrolytic oil, non-condensed gases, and biochar, as well as the percentages and treatment methods. I am ensuring that the yields and processing pathways are clear, including the initial crude pyrolytic oil's conversion to distillation and refinement. Furthermore, I've corrected a content error in the product streams table.


**Calculating Product Yields**

I'
<truncated 1061 bytes>
is a crucial step to determine the total net output of the process. I am making sure the final value is in line with the figures included in the project files.


**Evaluating Oil Quality Data**

I'm now evaluating the real-world pyrolysis oil analysis data. I've noted the average calorific value of approximately 41.0 MJ/kg, with a range between 36.575 and 43.405 MJ/kg, comparable to commercial diesel. The ash and impurity content is low. However, I'm focusing on the sulfur content (0.19%, exceeding the 0.05% limit), the flash point (below 40°C), and the cetane index as critical points for refining.


**Assessing Cetane Improvements**

I'm now assessing how to improve the cetane number (CN) of the oil. The initial CN of approximately 35.6 is below the minimum required for direct diesel engine use. I am internalizing the need for chemical refining and additive blending. I am now concentrating on improving the cetane index and flash point before integration into the generator.


**Analyzing Economic Significance**

I'm now analyzing the economic significance of the refined oil yields. The value aligns with the summarized project data. Non-condensed gas is an important internal energy flow, and carbon char is a solid stream that needs re-evaluation for its composition and environmental risks before utilization. Converting the oil yield at the estimated commercial oil price of about 22,000 VND/L, it results in a significant economic value, serving as a basis for efficiency analysis within the engineering project.


**Refining Analysis Data**

I'm now refining my analysis of the pyrolysis oil data, specifically its performance against QCVN 01:2022/BKHCN (diesel fuel standards). I've confirmed that the sulfur content (0.19%) exceeds the DO 0.05S limit, and the flash point is below the 55°C requirement. Also, the cetane number (CN) of about 35.6 is below the minimum required 46 for diesel engines. This solidifies the need for chemical refining and additive blending for cetane improvement before use in the generator.


","tool_calls":[{"name":"multi_replace_file_content","args":{"Description":""Fix bad replacements in Chuong 3"","Instruction":""Sửa lỗi do fuzzy matching lần trước gán nhầm dòng. Cập nhật lại các đoạn cân bằng sản phẩm, chưng cất và đánh giá chất lượng dầu vào đúng vị trí."","ReplacementChunks":"[{"AllowMultiple":false,"EndLine":364,"ReplacementContent":"\\midrule\
Dầu nhiệt phân thô & 45 & 20,74 & Chuyển sang chưng cất và tinh chế \\\\\
\\midrule\
Khí không ngưng & 30 & 13,82 & Hồi lưu làm nhiên liệu gia nhiệt lò \\\\\
\\midrule\
Than carbon/biochar & 25 & 11,52 & Kiểm nghiệm để tận dụng hoặc xử lý an toàn \\\\\
\\midrule\
Tổng cộng & 100 & 46,08 & \\\\\
\\midrule\
\\bottomrule\
\\end{tabular}\
\\end{table}\
\
Kết quả cân bằng cho thấy từ 46,08 tấn RDF/ngày, hệ thống nhiệt phân tạo ra 20,74 tấn dầu thô, 13,82 tấn khí không ngưng và 11,52 tấn than carbon. Trong ba dòng sản phẩm, dầu thô là dòng có giá trị kinh tế cao nhất nhưng cũng là dòng yêu cầu kiểm soát chất lượng nghiêm ngặt nhất trước khi đưa sang khâu chưng cất và tinh chế. Khí không ngưng là dòng năng lượng nội bộ quan trọng, còn than carbon là dòng rắn cần được đánh giá lại về thành phần và nguy cơ môi trường trước khi tận dụng.","StartLine":341,"TargetContent":"\\midrule\
Dầu nhiệt phân thô & 45 & 20,74 & Chuyển sang chưng cất và tinh chế \\\\\
\\midrule\
Khí không ngưng & 30 & 13,82 & Hỗn hợp dầu trung bình và dầu nhẹ ngưng tụ sẽ được đưa vào hệ thống chưng cất phụ trợ để tách phân đoạn. Từ tổng khối lượng dầu thô là 20,74 tấn/ngày, trừ đi 5,18 tấn dầu nặng (có thể dùng trộn nhựa đường hoặc đốt nồi hơi), còn lại 15,56 tấn dầu được chưng cất. Trong đó khoảng 5--10\\% tiếp tục tạo cặn (pitch) và 90--95\\% bay hơi thành phân đoạn xăng/diesel. \
\
Phân đoạn dầu chưng cất tiếp tục qua quá trình tinh chế hóa học bao gồm xử lý bằng axit sulfuric đậm đặc (2--5\\%), rửa kiềm bằng dung dịch \\ce{NaOH} 3\\% và hấp phụ bằng đất sét hoạt hóa. Định mức tiêu hao hóa chất dự ki
<truncated 9408 bytes>","TargetFile":""/Users/minhz/Desktop/HUST/KS/Đồ án tốt nghiệp kỹ sư/Engineer Thesis/Chuong/3_Chuong_3.tex"","toolAction":""Editing file"","toolSummary":""Update file contents""}}]}
{"step_index":33,"source":"MODEL","type":"CODE_ACTION","status":"DONE","created_at":"2026-06-10T07:42:06Z","content":"Created At: 2026-06-10T07:42:06Z
Completed At: 2026-06-10T07:42:06Z
The following changes were made by the multi_replace_file_content tool to: /Users/minhz/Desktop/HUST/KS/Đồ án tốt nghiệp kỹ sư/Engineer Thesis/Chuong/3_Chuong_3.tex. If relevant, proactively run terminal commands to execute this code for the USER. Don't ask for permission.
[diff_block_start]
@@ -341,11 +341,7 @@
 \midrule
 Dầu nhiệt phân thô & 45 & 20,74 & Chuyển sang chưng cất và tinh chế \\
 \midrule
-Khí không ngưng & 30 & 13,82 & Hỗn hợp dầu trung bình và dầu nhẹ ngưng tụ sẽ được đưa vào hệ thống chưng cất phụ trợ để tách phân đoạn. Từ tổng khối lượng dầu thô là 20,74 tấn/ngày, trừ đi 5,18 tấn dầu nặng (có thể dùng trộn nhựa đường hoặc đốt nồi hơi), còn lại 15,56 tấn dầu được chưng cất. Trong đó khoảng 5--10\% tiếp tục tạo cặn (pitch) và 90--95\% bay hơi thành phân đoạn xăng/diesel. 
-
-Phân đoạn dầu chưng cất tiếp tục qua quá trình tinh chế hóa học bao gồm xử lý bằng axit sulfuric đậm đặc (2--5\%), rửa kiềm bằng dung dịch \ce{NaOH} 3\% và hấp phụ bằng đất sét hoạt hóa. Định mức tiêu hao hóa chất dự kiến cho mỗi tấn dầu thô là: axit sulfuric 16--40 kg, xút \ce{NaOH} 24 kg, và đất sét hoạt tính 16--40 kg. Quá trình này giúp loại bỏ hợp chất chứa lưu huỳnh, nhựa không bão hòa và cải thiện độ màu.
-
-Sau các bước rửa nước, thải bỏ pha axit và mất mát trong vật liệu hấp phụ, lượng dầu nhẹ tinh chế giảm nhẹ so với khối lượng ban đầu. Với giả thiết tổn thất khoảng 0,6\%, sản lượng dầu thành phẩm cuối cùng được tính như sau: \\
+Khí không ngưng & 30 & 13,82 & Hồi lưu làm nhiên liệu gia nhiệt lò \\
 \midrule
 Than carbon/biochar & 25 & 11,52 & Kiểm nghiệm để tận dụng hoặc xử l
<truncated 4406 bytes>
ể. Tuy nhiên, trong phạm vi đồ án kỹ sư, kết quả này chỉ được sử dụng như cơ sở tham chiếu cho phân tích hiệu quả chứ chưa phải là dự báo doanh thu thương mại chính thức.
 
+\subsection{Phân tích chất lượng dầu nhiệt phân thực tế}
+
+Kết quả phân tích các mẫu dầu nhiệt phân thực tế tại nhà máy cho thấy dầu có nhiệt trị trung bình đạt khoảng 41,0 MJ/kg (kết quả kiểm định dao động từ 36,575 đến 43,405 MJ/kg tùy mẫu), tương đương nhiên liệu diesel thương mại (42--45 MJ/kg). Hàm lượng tro và tạp chất rất thấp (	extless{}0,01\%), khối lượng riêng ở 15$^\circ$C khoảng 0,779--0,826 kg/L. 
+
+Tuy nhiên, phân tích thực tế cũng chỉ ra dầu nhiệt phân thô tồn tại một số hạn chế so với QCVN 01:2022/BKHCN (áp dụng cho nhiên liệu điêzen). Cụ thể, hàm lượng lưu huỳnh trung bình 0,19\% (vượt ngưỡng 0,05\% của DO 0,05S), điểm chớp cháy cốc kín 	extless{}40$^\circ$C (thấp hơn yêu cầu tối thiểu 55$^\circ$C), và chỉ số cetane (CN) đạt khoảng 35,6 (thấp hơn mức tối thiểu 46). Chỉ số cetane thấp làm kéo dài thời gian trễ cháy, giảm hiệu suất và tăng phát thải khi dùng trực tiếp cho động cơ diesel. Vì vậy, công đoạn tinh chế hóa học và phối trộn phụ gia là bắt buộc để cải thiện chỉ số cetane và điểm chớp cháy trước khi đưa vào cụm máy phát điện.
+
 \subsection{Đặc điểm và hướng xử lý than carbon sau nhiệt phân}
 
 Than carbon (carbon black hay pyrolysis char) là phần rắn còn lại sau nhiệt phân, gồm các thành phần chính:
[diff_block_end]

Please note that the above snippet only shows the MODIFIED lines from the last change. It shows up to 3 lines of unchanged lines before and after the modified lines. The actual file contents may have many more lines not shown."}
{"step_index":34,"source":"SYSTEM","type":"EPHEMERAL_MESSAGE","status":"DONE","created_at":"2026-06-10T07:42:06Z","content":"The following is an <EPHEMERAL_MESSAGE> not actually sent by the user. It is provided by the system as a set of reminders and general important information to pay attention to. Do NOT respond to this message, just act accordingly.

<EPHEMERAL_MESSAGE>
<planning_mode>
You're in planning mode. Exercise judgement on whether the user's last request warrants a plan and take action accordingly.
If the user's request warrants a plan then follow this workflow:
1. Research first before taking actions. Do not use your code editing tools or run modifying commands during this time.
2. Create an implementation plan artifact with `request_feedback = true` to request approval. Include any open questions for the user directly in the implementation plan. Don't use the ask_question tool to ask these questions. Stop calling tools and allow the user to review the plan and provide feedback or approval.
3. After receiving approval, begin execution. Create and update the task artifact as you work.
4. Verify any work that you have done.
5. Create a walkthrough artifact to present your work to the user. When writing artifacts, make sure to follow the artifact formatting guidelines and use github alerts to highlight critical sections.
If the user's request does NOT warrant a plan then continue your work WITHOUT making a plan or requesting user feedback.
</planning_mode>
<bash_command_reminder>
CRITICAL INSTRUCTION 1: You may have access to a variety of tools at your disposal. Some tools may be for a specific task such as 'view_file' (for viewing contents of a file). Others may be very broadly applicable such as the ability to run a command on a terminal. Always prioritize using the most specific tool you can for the task at hand. Here are some rules: (a) NEVER run cat inside a bash command to create a new file or append to an existing file. (b) ALWAYS use grep_search instead of running grep inside a bash command unless absolutely needed. (c) DO NOT use ls for listing, cat for viewing, grep for finding, sed for replacing. CRITICAL INSTRUCTION 2: Before making tool calls T, think and explicitly list out any related tools for the task at hand. You can only execute a set of tools T if all other tools in the list are either more generic or cannot be used for the task at hand. ALWAYS START your thought with recalling critical instructions 1 and 2. In particular, the format for the start of your thought block must be '...94>thought
CRITICAL INSTRUCTION 1: ...
CRITICAL INSTRUCTION 2: ...'.
</bash_command_reminder>
</EPHEMERAL_MESSAGE>"}
{"step_index":35,"source":"MODEL","type":"PLANNER_RESPONSE","status":"DONE","created_at":"2026-06-10T07:42:06Z","thinking":"**Prioritizing Tool Specificity**

I'm focusing now on tool specificity. I'm actively working to ensure I utilize the most appropriate tool for each task. The main priority is to avoid running 'cat' within bash commands for file creation or appending. Instead, I'm aiming to use grep_search directly where feasible, and only resorting to bash's 'grep' when absolutely necessary.


**Refining Tool Usage**

I'm now prioritizing tool usage and task breakdowns. I'm ensuring I select the best tool for the job. Specifically, for listing, viewing, searching, and replacing, I will use `ls`, `cat`, `grep`, and `sed`, respectively. Before executing any toolset, I will explicitly list all relevant tools, executing only when all others are less suitable. I've successfully updated `2_Chuong_2.tex` and `3_Chuong_3.tex` with content from `NCKH2026.tex` and am checking the context to avoid unprompted compiling. I am also identifying and preparing to incorporate references to tables within `NCKH2026.tex`.


","tool_calls":[{"name":"view_file","args":{"AbsolutePath":""/Users/minhz/Desktop/HUST/KS/Đồ án tốt nghiệp kỹ sư/Engineer Thesis/Chuong/1_Chuong_1.tex"","EndLine":"200","StartLine":"1","toolAction":""Viewing file"","toolSummary":""View Chapter 1""}}]}
{"step_index":36,"source":"MODEL","type":"VIEW_FILE","status":"DONE","created_at":"2026-06-10T07:42:14Z","content":"Created At: 2026-06-10T07:42:14Z
Completed At: 2026-06-10T07:42:14Z
File Path: `file:///Users/minhz/Desktop/HUST/KS/%C4%90o%CC%82%CC%80%20a%CC%81n%20to%CC%82%CC%81t%20nghie%CC%A3%CC%82p%20ky%CC%83%20su%CC%9B/Engineer%20Thesis/Chuong/1_Chuong_1.tex`
Total Lines: 572
Total Bytes: 62822
Showing lines 1 to 200
The following code has been modified to include a line number before every line, in the format: <line_number>: <original_line>. Please note that any changes targeting the original code should remove the line number, colon, and leading space.
1: % !TeX root = ../main.tex
2: 
3: \documentclass[../main.tex]{subfiles}
4: \begin{document}
5: 
6: \definecolor{OrganicGreen}{HTML}{2D5F2E}  % Xanh lá đậm - Rác hữu cơ
7: \definecolor{PlasticOrange}{HTML}{E67E22} % Cam đất - Nhựa
8: \definecolor{SandBrown}{HTML}{A0522D}    % Nâu đất - Đất cát
9: \definecolor{PaperBlue}{HTML}{2980B9}     % Xanh dương - Giấy
10: \definecolor{GlassTeal}{HTML}{1ABC9C}     % Xanh ngọc - Thủy tinh
11: \definecolor{MetalGray}{HTML}{95A5A6}     % Xám chì - Kim loại
12: \definecolor{OtherGray}{HTML}{D3D3D3}     % Xám nhạt - Khác
13: 
14: \chapter{TỔNG QUAN}
15: 
16: \section{Hiện trạng quản lý và xử lý chất thải rắn sinh hoạt tại Việt Nam}
17: 
18: \subsection{Đặc điểm và thành phần rác thải sinh hoạt}
19: 
20: Chất thải rắn sinh hoạt (CTRSH) tại Việt Nam mang những đặc điểm riêng biệt, phản ánh rõ nét trình độ phát triển kinh tế - xã hội và văn hóa sinh hoạt của người dân. Thành phần CTRSH nổi bật với tỷ lệ chất hữu cơ cao, thường dao động từ 50$-$70$\%$ tổng khối lượng rác. Các chất hữu cơ này chủ yếu bao gồm thức ăn thừa, rau quả, lá cây và các vật liệu dễ phân hủy sinh học khác. Đây là đặc điểm điển hình của các quốc gia đang phát triển, nơi thói quen tiêu dùng và chế biến thực phẩm tại hộ gia đình vẫn còn
<truncated 17919 bytes>
ợng hữu ích. Tùy thuộc vào tác nhân oxy hóa và nhiệt độ phản ứng, các công nghệ này được chia thành ba nhánh chính.
193: 
194: 	extbf{Đốt rác có thu hồi năng lượng (Waste-to-Energy - WtE):} Rác thải được đốt cháy hoàn toàn trong môi trường dư oxy ở nhiệt độ từ 850$-$1100$^\circ$C. Nhiệt lượng sinh ra làm hóa hơi nước trong lò hơi, kéo tuabin phát điện. Đây là công nghệ trưởng thành nhất. Đối với rác độ ẩm cao như ở Việt Nam, lò đốt thường cần sấy sơ bộ hoặc phun thêm dầu mồi để ổn định quá trình cháy.
195: 
196: 	extbf{Khí hóa (Gasification):} Quá trình chuyển hóa chất thải chứa carbon trong điều kiện thiếu oxy ở nhiệt độ cao (800$-$1200$^\circ$C), tạo ra khí tổng hợp (syngas) chứa \ce{CO}, \ce{H2}, \ce{CH4} \cite{Arena2012}. Nhược điểm lớn của khí hóa là đòi hỏi nguyên liệu (RDF) phải rất đồng nhất, độ ẩm thấp. Dòng khí tổng hợp chứa hắc ín (tar), bụi và các chất ăn mòn, làm cho hệ thống làm sạch khí trở nên phức tạp và đắt đỏ.
197: 
198: 	extbf{Nhiệt phân (Pyrolysis):} Quá trình phân hủy nhiệt chất hữu cơ trong môi trường không có oxy, ở nhiệt độ 350$-$600$^\circ$C đối với rác nhựa và RDF \cite{Bridgwater2012,Sharuddin2016}. Sản phẩm tạo ra gồm dầu nhiệt phân (pha lỏng), khí không ngưng (pha khí), và than carbon (pha rắn). Khác với khí hóa, nhiệt phân sinh ra nhiên liệu lỏng có thể lưu trữ, chưng cất thành nhiên liệu tương đương diesel, rất phù hợp với thành phần rác chứa nhiều bao bì nhựa, nylon \cite{Miandad2016,Williams2007}.
199: 
200: \begin{table}[htbp]
The above content does NOT show the entire file contents. If you need to view any lines of the file which were not shown to complete your task, call this tool again to view those lines.
"}
{"step_index":38,"source":"SYSTEM","type":"EPHEMERAL_MESSAGE","status":"DONE","created_at":"2026-06-10T07:42:14Z","content":"The following is an <EPHEMERAL_MESSAGE> not actually sent by the user. It is provided by the system as a set of reminders and general important information to pay attention to. Do NOT respond to this message, just act accordingly.

<EPHEMERAL_MESSAGE>
<planning_mode>
You're in planning mode. Exercise judgement on whether the user's last request warrants a plan and take action accordingly.
If the user's request warrants a plan then follow this workflow:
1. Research first before taking actions. Do not use your code editing tools or run modifying commands during this time.
2. Create an implementation plan artifact with `request_feedback = true` to request approval. Include any open questions for the user directly in the implementation plan. Don't use the ask_question tool to ask these questions. Stop calling tools and allow the user to review the plan and provide feedback or approval.
3. After receiving approval, begin execution. Create and update the task artifact as you work.
4. Verify any work that you have done.
5. Create a walkthrough artifact to present your work to the user. When writing artifacts, make sure to follow the artifact formatting guidelines and use github alerts to highlight critical sections.
If the user's request does NOT warrant a plan then continue your work WITHOUT making a plan or requesting user feedback.
</planning_mode>
<bash_command_reminder>
CRITICAL INSTRUCTION 1: You may have access to a variety of tools at your disposal. Some tools may be for a specific task such as 'view_file' (for viewing contents of a file). Others may be very broadly applicable such as the ability to run a command on a terminal. Always prioritize using the most specific tool you can for the task at hand. Here are some rules: (a) NEVER run cat inside a bash command to create a new file or append to an existing file. (b) ALWAYS use grep_search instead of running grep inside a bash command unless absolutely needed. (c) DO NOT use ls for listing, cat for viewing, grep for finding, sed for replacing. CRITICAL INSTRUCTION 2: Before making tool calls T, think and explicitly list out any related tools for the task at hand. You can only execute a set of tools T if all other tools in the list are either more generic or cannot be used for the task at hand. ALWAYS START your thought with recalling critical instructions 1 and 2. In particular, the format for the start of your thought block must be '...94>thought
CRITICAL INSTRUCTION 1: ...
CRITICAL INSTRUCTION 2: ...'.
</bash_command_reminder>
</EPHEMERAL_MESSAGE>"}
{"step_index":39,"source":"MODEL","type":"PLANNER_RESPONSE","status":"DONE","created_at":"2026-06-10T07:42:14Z","thinking":"**Prioritizing Tool Specificity**

I'm focusing now on tool specificity. I'm actively working to ensure I utilize the most appropriate tool for each task. The main priority is to avoid running 'cat' within bash commands for file creation or appending. Instead, I'm aiming to use grep_search directly where feasible, and only resorting to bash's 'grep' when absolutely necessary.


","tool_calls":[{"name":"view_file","args":{"AbsolutePath":""/Users/minhz/Desktop/HUST/KS/Đồ án tốt nghiệp kỹ sư/Engineer Thesis/Chuong/1_Chuong_1.tex"","EndLine":"400","StartLine":"200","toolAction":""Viewing file"","toolSummary":""View Chapter 1""}}]}
{"step_index":40,"source":"MODEL","type":"VIEW_FILE","status":"DONE","created_at":"2026-06-10T07:42:18Z","content":"Created At: 2026-06-10T07:42:18Z
Completed At: 2026-06-10T07:42:18Z
File Path: `file:///Users/minhz/Desktop/HUST/KS/%C4%90o%CC%82%CC%80%20a%CC%81n%20to%CC%82%CC%81t%20nghie%CC%A3%CC%82p%20ky%CC%83%20su%CC%9B/Engineer%20Thesis/Chuong/1_Chuong_1.tex`
Total Lines: 572
Total Bytes: 62822
Showing lines 200 to 400
The following code has been modified to include a line number before every line, in the format: <line_number>: <original_line>. Please note that any changes targeting the original code should remove the line number, colon, and leading space.
200: \begin{table}[htbp]
201: \centering
202: \caption{So sánh khái quát một số phương pháp xử lý CTRSH}
203: \label{tab:so-sanh-phuong-phap-xu-ly}
204: \setlength{	abcolsep}{5pt}
205: \scalebox{0.9}{
206: \begin{tabular}{>{\raggedright\arraybackslash}p{0.18	extwidth}p{0.26	extwidth}p{0.32	extwidth}>{\raggedright\arraybackslash}p{0.20	extwidth}}
207: 	oprule
208: 	extbf{Phương pháp} & 	extbf{Ưu điểm} & 	extbf{Hạn chế} & 	extbf{Sản phẩm chính} \\
209: \midrule
210: Chôn lấp hợp vệ sinh & Đơn giản, tiếp nhận nhiều loại rác & Chiếm đất, phát sinh nước rỉ rác và khí nhà kính & Khí bãi rác, nước rỉ rác \\
211: Ủ compost & Phù hợp rác hữu cơ đã phân loại & Thị trường sản phẩm hạn chế, không xử lý rác nhựa & Phân hữu cơ \\
212: Đốt trực tiếp (WtE) & Giảm nhanh thể tích rác, thu hồi năng lượng lớn & Yêu cầu kiểm soát khí thải nghiêm ngặt, bất lợi với rác ẩm & Điện năng, nhiệt, tro xỉ \\
213: Khí hóa & Thu hồi khí nhiên liệu, hiệu suất lý thuyết cao & Yêu cầu RDF đồng nhất, làm sạch khí phức tạp & Khí tổng hợp, tro xỉ \\
214: Nhiệt phân & Tạo dầu lỏng dễ lưu trữ, giảm bụi do không đốt trực tiếp & Cần phân loại và tinh chế sản phẩm dầu & Dầu nhiệt phân, khí, than carbon \\
215: \bottomrule
216: \end{tabular}%
217: }
218: \end{tabl
<truncated 20975 bytes>
 62,56\% khối lượng; nhựa và bao nylon chiếm 10,16\%; giấy chiếm 7,42\%; đất cát chiếm 7,76\%; thủy tinh chiếm 2,28\% và kim loại chiếm 1,45\%. Về tính chất cháy, nhóm chất thải cháy được chiếm khoảng 86,62\%, còn nhóm không cháy chiếm khoảng 13,38\%. Độ ẩm trung bình của mẫu rác là 56,3\%, độ tro khoảng 8,5\% \cite{BaocaoHoiAn2024}.
358: 
359: \begin{figure}[H]
360: \centering
361: \makeatletter
362: \def\pgfpie@legend#1{%
363:   \coordinate[xshift=0.8cm,
364:   yshift={(	he\pgfpie@sliceLength*0.5+1)*0.8cm}] (pgfpie@legendpos) at
365:   (current bounding box.east);
366: 
367:   \scope[node distance=0.8cm]
368:     \foreach \pgfpie@p/\pgfpie@t [count=\pgfpie@i from 0] in {#1}
369:     {
370:       \pgfpie@findColor{\pgfpie@i}
371:       \
ode[draw, fill={\pgfpie@thecolor}, \pgfpie@style, below of={pgfpie@legendpos}, label={0:{\pgfpie@t}}] (pgfpie@legendpos) {};
372:     }
373:   \endscope
374: }
375: \makeatother
376: \scalebox{0.9}{
377: \begin{tikzpicture}
378:     \pie[
379:         style=drop shadow,
380:         hide number,
381:         text=legend,
382:         radius=3.5,
383:         color={OrganicGreen, PlasticOrange, SandBrown, PaperBlue, GlassTeal, MetalGray, OtherGray},
384:         explode=0.1
385:     ]{
386:         62.56/{Rác hữu cơ (62,56\%)},
387:         10.16/{Nhựa \& bao nylon (10,16\%)},
388:         7.76/{Đất cát (7,76\%)},
389:         7.42/{Giấy (7,42\%)},
390:         2.28/{Thủy tinh (2,28\%)},
391:         1.45/{Kim loại (1,45\%)},
392:         8.37/{Thành phần khác (8,37\%)}
393:     }
394: \end{tikzpicture}
395: }
396: \caption{Thành phần khối lượng CTRSH tại Hội An}
397: \label{fig:thanh-phan-rac}
398: \end{figure}